\begin{frame}{어텐션 출력의 볼록성: 볼록 껍질 안에만 존재}
  마지막 단계(가중합)가 \textbf{가중치의 합이 1인 가중 평균}이라는
  사실은 출력의 \emph{범위}에 강한 제약을 준다:
  \begin{itemize}
    \item 각 출력 $o_i = \sum_j w_{ij} v_j$ ($\sum_j w_{ij} = 1$,
      $w_{ij} \ge 0$)은 \textbf{$V$ 의 원행들의 볼록 결합}
      $\Rightarrow$ $V$ 원행들로 만드는 \textbf{볼록 껍질
      (convex hull) 안에 항상 존재}.
    \item \textbf{네 단어 예 숫자로 확인}: ``동물''의 출력
      $[2.898, 0.133, 0.010]$ — $V$ 원행(= $E$ 의 원행)들의
      범위가 각 차원 0$\sim$3인데, 각 성분은 모두 0$\sim$3 안에
      있고 \textbf{성분 합도 3.041로 $V$ 원행의 성분 합 (= 3)과
      거의 같}다. ``균등 어텐션''(모두 1/4)의 출력은 단순 평균
      $[1.45, 0.825, 0.75]$ 로, 실제 출력과는 \emph{완전히
      다른} 벡터 — 가중치가 ``분포''에 따라 출력은 볼록 껍질
      안에서 \emph{아무 위치나} 만들 수 있다.
  \end{itemize}
  \vskip0.4em
  \begin{block}{실전에서 이어지는 두 가지}
    (1) \textbf{출력의 스케일이 입력 스케일의 범위 안에 갇혀
    있다} — $V$ 가 큰 값으로 학습되면 출력도 그 범위 안에서만
    ``미세 조정'' (때로는 \emph{장점 — 안정성}, 때로는
    \emph{한계} — ``$V$ 가 표현할 수 없는 값''은 한 층으로
    만들 수 없음). (2) \textbf{가중치가 극단(one-hot)이면 출력은
    $V$ 의 어떤 \emph{원행}에, 균등이면 \emph{평균}에} 가까워진다
    — softmax는 ``어떤 원행을 얼마나 섞을지''의 \textbf{비율},
    \emph{새로운} 값이 아니라 \emph{섞임}을 만드는 도구.
  \end{block}
  \vskip0.3em
  {\footnotesize \kq{주의: ``합이 1인 가중치''가 보장하는 것은
  ``출력이 $V$ 원행들 \emph{사이에 있다}(볼록성)''는 것이지,
  출력의 \emph{스케일}이 $V$ 의 어떤 스케일을 \emph{정확히}
  유지한다는 보장은 \emph{아니다} (V 원행끼리 스케일이 다르면
  출력 스케일은 그 스케일들의 \emph{가중 평균}이 된다).}}
\end{frame}
